% !TEX root = ../main.tex
\clearpage
\section{Lecciones de despliegues y propuesta de piloto}
\label{sec:piloto}
\subsection{Resultados y límites de las experiencias previas}
\textbf{El alivio documental no implica necesariamente consultas más breves.} En Suecia se comunica ahorro de tiempo percibido, mientras que en Quirónsalud las mejoras de experiencia coexisten con una duración media de consulta ligeramente mayor entre los usos de Scribe. El ensayo de UCLA muestra además que el resultado puede diferir entre productos~\cite{C7S8N9NA,U5GHQW86,PR93N8FN}. Para el hospital, la consecuencia es evaluar por separado carga percibida, tiempo total de documentación y calidad clínica. La experiencia de Cleveland Clinic subraya también el papel de la formación y el soporte en la adopción~\cite{F5FVUMMX}. Permanente añade una advertencia: el uso se concentra en una parte de los profesionales~\cite{tierney2025}. Los casos siguientes orientan esas decisiones; no constituyen una comparación directa entre proveedores.

\begin{small}
\begin{longtable}{L{.23\linewidth}L{.69\linewidth}}
\toprule
Experiencia & Resultado y consecuencia para el proyecto \\
\midrule\endhead
UCLA, Estados Unidos & Ensayo aleatorizado con 238 médicos: reducción del tiempo de documentación con Nabla, sin diferencia significativa con DAX. Cerca del 15\% de los médicos asignados a escribas no los utilizó. Conviene incluir a todos los participantes y registrar el abandono~\cite{PR93N8FN}. \\
The Permanente Medical Group, Estados Unidos & Seguimiento de 63 semanas: 7.260 médicos y 2,58 millones de encuentros; ahorro agregado anual estimado de 15.791 horas, comparando cambios frente a no usuarios. Diseño observacional, sin asignación aleatoria. La escala permite un ahorro acumulado importante, pero no demuestra un beneficio uniforme. Conviene analizar utilización y carga inicial por profesional~\cite{tierney2025}. \\
Abridge, seis sistemas de Estados Unidos & Estudio antes/después de 263 profesionales: mejoras autodeclaradas de bienestar y carga documental tras 30 días. La participación voluntaria y la ausencia de un control contemporáneo limitan la atribución causal. Conviene medir bienestar sin equipararlo a ahorro objetivo~\cite{olson2025}. \\
Cleveland Clinic, Estados Unidos & Informe de implantación de Ambience: más de 4.000 profesionales incorporados en cuatro meses, con gobernanza, formación y soporte. La coautoría del proveedor y el diseño descriptivo limitan la interpretación clínica. El caso aporta un modelo organizativo~\cite{F5FVUMMX}. \\
Tandem, Suecia & Estudio observacional de 236.153 notas; encuesta de 177 usuarios con utilización sostenida. Reducción autodeclarada del 29\% en tiempo por nota. La selección de usuarios y la edición fuera de la aplicación pueden sesgar la estimación del beneficio~\cite{C7S8N9NA}. \\
Quirónsalud, España & Más de 2,33 millones de usos ambulatorios, con mejoras percibidas; duración media de 15,01 minutos con Scribe y 14,65 sin él. El acuerdo semántico de 87,4--89,2\% compara borrador y texto editado, sin establecer exactitud clínica independiente. Satisfacción y duración requieren análisis separado~\cite{U5GHQW86}. \\
GOSH y otros centros NHS, Londres & Evaluación de Tortus en nueve centros y más de 17.000 encuentros. El hospital comunica menor carga documental y mayor tiempo de interacción. Se trata de un estudio antes/después no aleatorizado, difundido como informe y preprint; los efectos no se extrapolan automáticamente al español~\cite{gosh2025,wray2025}. \\
\bottomrule
\end{longtable}
\end{small}
En España, el Hospital Fundación Alcorcón aporta además una referencia institucional de uso de Invox Genesis en Psiquiatría, descrita en la sección~\ref{sec:candidatos}. Esa adopción demuestra experiencia de implantación, pero no sustituye una evaluación comparativa del rendimiento.

\subsection{Secuencia de trabajo propuesta}
Se propone un proceso progresivo, sujeto a disponibilidad presupuestaria y a la definición del hospital. El apéndice~\ref{app:evaluacion-piloto} desarrolla opciones de diseño y medidas de resultado para su revisión. Es una base de discusión; no fija todavía el protocolo, la muestra ni los umbrales de aceptación.
\begin{enumerate}
\item \textbf{Preparación local.} Seleccionar servicios, profesionales interesados y condiciones técnicas; registrar la situación inicial de documentación y resolver el tratamiento de datos permitido.
\item \textbf{Conjunto común de consultas de prueba.} Valorar la recogida de audio y notas con finalidad y condiciones específicas, sin que ello autorice su envío a proveedores. Si no es viable utilizar consultas reales, pueden emplearse escenarios simulados~\cite{rgpd2016,NNWFT2JX}.
\item \textbf{Comparación previa al uso asistencial.} Probar entre uno y tres proveedores sobre casos comunes, según coste y requisitos, con revisión clínica de los borradores y sin incorporarlos a HCIS como documentación asistencial.
\item \textbf{Piloto asistencial limitado.} Evaluar los candidatos que superen la fase previa, con revisión profesional y alternativa ante fallos. La ampliación dependerá de criterios acordados antes del inicio.
\end{enumerate}

La asignación debería equilibrar especialidades, complejidad y experiencia digital; asignar un producto a cada servicio confundiría ambos efectos. Los profesionales interesados pueden facilitar el inicio, pero deberán registrarse también la falta de uso, los rechazos y el abandono~\cite{PR93N8FN,FKHM4UFL}.

El banco de consultas permite comparar notas; para estimar ahorro asistencial será necesario comparar flujos de trabajo. Las alternativas de diseño, la referencia clínica y las medidas propuestas se desarrollan en el apéndice~\ref{app:evaluacion-piloto}.
